\documentclass[10pt,a4paper]{amsart}
\usepackage[margin=19mm]{geometry}
\usepackage{amsmath,amssymb,mathtools}
\usepackage[scr=rsfs,cal=euler]{mathalfa}
\usepackage{xfrac}
\usepackage[unicode,hidelinks,bookmarks=false]{hyperref}
\newcommand{\ie}{i.e.\ }
\newcommand{\cf}{cf.\ }
\newcommand{\resp}{respectively\ }
\newcommand{\Subsection}{\subsection}
\providecommand{\nobreakdash}{\nobreak\hbox{-}}
\setlength{\emergencystretch}{2em}
\multlinegap0pt
\usepackage{bm}
\usepackage{bbm}
\usepackage{dsfont}
\usepackage{xspace}
\usepackage{cancel}
\usepackage{mathtools}
\usepackage{amsthm}
\makeatletter
\@ifundefined{theorem}{\newtheorem{theorem}{Theorem}}{}
\makeatother
\title[On the Graphical $r$-Stirling Numbers of the First Kind for ]{On the Graphical $r$-Stirling Numbers of the First Kind for Specific Graph Families}
\author{Daniel Yaqubi, Madjid Mirzavaziri}
\date{Source: arXiv:2602.02046v1 (CC BY 4.0)}
\begin{document}
\maketitle

\noindent\emph{Source and licence.} On the Graphical $r$-Stirling Numbers of the First Kind for Specific Graph Families. Version arXiv:2602.02046v1 (CC BY 4.0), distributed under the Creative Commons Attribution 4.0 International licence, as recorded in the arXiv licence metadata for this exact version and shown on the abstract page https://arxiv.org/abs/2602.02046v1. Full rights notice: licenses/arxiv-2602.02046-CC-BY-4.0.txt.\par\smallskip

\noindent\textit{Editorial scope.} Counted unit: the Theorem whose source label is \texttt{thm:cycle\_stirling} in the article (source0.tex lines 1055--1065, source label \texttt{thm:cycle\_stirling}), delivered together with the article's own complete original proof of it (source0.tex lines 1067--1111, one \texttt{proof} environment of 45 lines). No mathematics is written, repaired or re-derived: statement and proof are the article's own text line for line. Required context retained uncounted: none. Every definition, display and label of those lines is retained unchanged. Omitted from this package and not redistributed: 1--1054, 1066--1066, 1112--1999. Nothing inside a retained excerpt is omitted or altered: every retained body is delivered exactly as the article's lines are written, so no omission receipt of any kind accompanies this package. Citations: the retained excerpts carry no citation command at all, so no bibliography entry of the article is transcribed and no reference is redistributed. Reference adaptation: none was needed and none was made; every reference inside the delivered unit resolves inside it, so no unresolved reference or citation is produced. Printed slips kept literal: no printed word, symbol, hypothesis or number of the counted statement or of its proof was altered. Layout and class: the article's own class is \texttt{amsart} with packages including fixltx2e, amsmath, amsthm, amscd, amsfonts, amssymb, graphicx, color, float, pgf, tikz, fontenc, latexsym, wasysym; the wrapper is the lane's pinned \texttt{amsart} editorial wrapper at 10pt on A4, which restates the theorem-like environments the retained text needs and adds the article's own macro definitions that the retained text needs (none), with extra wrapper packages (bm, bbm, dsfont, xspace, cancel, mathtools). The delivered numbering is the unit's own. Assets: no retained span contains a figure environment, a graphics-inclusion command, a PDF-page inclusion, a table or a TikZ picture; no image file is copied into this package, the package declares no asset dependency and no image file is redistributed with it. Licence: version arXiv:2602.02046v1 of this article is distributed under the Creative Commons Attribution 4.0 International licence, as recorded in the arXiv licence metadata for this exact version and shown on the abstract page https://arxiv.org/abs/2602.02046v1. A full rights notice is delivered at licenses/arxiv-2602.02046-CC-BY-4.0.txt. Characters: every retained excerpt is pure ASCII, so the delivered engine, XeLaTeX with its default Latin Modern text font, reports no missing character.
\begin{theorem}\label{thm:cycle_stirling}
Let $n$ and $k$ be positive integers such that $k \geq r$. The graphical $r$-Stirling number of the first kind for the cycle graph $C_n$, denoted by ${C_n \brack k}_r$, is given by
\begin{equation*}
    {C_n \brack k}_r = \binom{k-r+2}{n-k}.
\end{equation*}
Furthermore, the cardinality of the set of all such $r$-graphical cycle partitions is 
\begin{equation*}
   B_r(C_n)=\sum_{k=r}^n {C_n \brack k}_r = F_{n-r+3},
\end{equation*}
where $F_m$ denotes the $m$-th Fibonacci number.
\end{theorem}

\begin{proof}
Let the vertex set of $C_n$ be $V = \{1, 2, \dots, n\}$ with the edge set $E = \{(i, i+1) : 1 \le i < n\} \cup \{(n, 1)\}$. An $r$-graphical cycle partition requires that the vertices $\{1, 2, \dots, r\}$ reside in distinct blocks, and each block must induce a cycle in $C_n$. Since the only cycles in $C_n$ are $1$-cycles (singletons) and $2$-cycles (edges), each block is either a singleton or an adjacent pair. 

We partition the collection of all valid cycle partitions into four exhaustive and mutually exclusive cases based on the neighborhood of the restricted boundary vertices $1$ and $r$:

\begin{itemize}
    \item \textbf{Case I:} Vertices $1$ and $r$ are both $1$-cycles. 
    The remaining $n-r$ vertices must be partitioned into $k-r$ cycles. This is equivalent to an $r$-restricted partition of a path graph $P_{n-r}$ into $k-r$ blocks. Thus, the number of such partitions is:
    \begin{equation*}
        {P_{n-r} \brack k-r} = \binom{k-r}{n-r-(k-r)} = \binom{k-r}{n-k}.
    \end{equation*}

    \item \textbf{Case II:} Vertex $1$ is a $1$-cycle, while vertex $r$ belongs to a $2$-cycle $\{r, r+1\}$. 
    The remaining $n-r-1$ vertices are partitioned into $k-r$ cycles on the path $P_{n-r-1}$, yielding:
    \begin{equation*}
        {P_{n-r-1} \brack k-r} = \binom{k-r}{n-r-1-(k-r)} = \binom{k-r}{n-k-1}.
    \end{equation*}

    \item \textbf{Case III:} Vertex $r$ is a $1$-cycle, while vertex $1$ belongs to a $2$-cycle $\{1, n\}$. 
    By symmetry with Case II, the number of ways to partition the remaining vertices into $k-r$ blocks on the path $P_{n-r-1}$ is:
    \begin{equation*}
        {P_{n-r-1} \brack k-r} = \binom{k-r}{n-k-1}.
    \end{equation*}

    \item \textbf{Case IV:} Both vertices $1$ and $r$ belong to $2$-cycles, specifically $\{1, n\}$ and $\{r, r+1\}$. 
    The remaining $n-r-2$ vertices are partitioned into $k-r$ blocks on the path $P_{n-r-2}$, giving:
    \begin{equation*}
        {P_{n-r-2} \brack k-r} = \binom{k-r}{n-r-2-(k-r)} = \binom{k-r}{n-k-2}.
    \end{equation*}
\end{itemize}

By summing the contributions from these four cases and repeatedly applying Pascal's Identity, $\binom{n}{k} = \binom{n-1}{k} + \binom{n-1}{k-1}$, we obtain:
\begin{align*}
    {C_n \brack k}_r &= \binom{k-r}{n-k} + 2\binom{k-r}{n-k-1} + \binom{k-r}{n-k-2} \\
    &= \left[ \binom{k-r}{n-k} + \binom{k-r}{n-k-1} \right] + \left[ \binom{k-r}{n-k-1} + \binom{k-r}{n-k-2} \right] \\
    &= \binom{k-r+1}{n-k} + \binom{k-r+1}{n-k-1} \\
    &= \binom{k-r+2}{n-k}.
\end{align*}

To compute the total number of partitions, we sum over the possible number of blocks $k$:
\begin{equation*}
      B_r(C_n)= \sum_{k=r}^n {C_n \brack k}_r = \sum_{k=r}^n \binom{k-r+2}{n-k}.
\end{equation*}
Letting $j = k-r$, the summation transforms to $\sum_{j=0}^{n-r} \binom{j+2}{n-r-j}$. Utilizing the identity for the sum of diagonals in Pascal's triangle, $\sum_{i=0}^{\lfloor m/2 \rfloor} \binom{m-i}{i} = F_{m+1}$, and setting $m = n-r+2$, we conclude that the sum equals $F_{n-r+3}$.
\end{proof}

\end{document}
